\section*{Limitations}
\paragraph{Domain and population coverage.}
Each domain is represented by a particular benchmark and prompt format. ACPBench questions, for example, are not equivalent to unrestricted plan construction. The atlas uses one generation seed per prompt/model and predominantly completed traces; missing annotations limit generalization. Shared families and task collections mean the seven checkpoints are not independent population samples. We infer neither scaling laws nor causal effects of training lineage. Multiple benchmarks per domain, instruction paraphrases, and repeated generations would help separate task demand, format, and sampling variation.

\paragraph{Measurement validity and human comparisons.}
The primary labels are generated by Gemma4-31B-it. Prior human validation of the source taxonomies supports the definitions, but does not transfer automatically to new judges, sentence-level decisions, and domains. The three-family audit preserves selected domain contrasts but shows substantial sentence-level disagreement, especially with Llama-3.1-8B-Instruct. This supports limited aggregate robustness, while correlated errors and semantic validity remain unresolved. Selected examples illustrate labels rather than validating them. Labels may track discourse style, and relative-position trajectories can obscure duration and termination. Event-aligned analyses and coding of productive revision versus repetition would strengthen measurement. We collect no human behavioral or neural data; cognitive-science connections concern task organization, not corresponding brain systems.

\paragraph{Intervention selection and inference.}
The main intervention experiments use Qwen3.5-9B and one inverse-readout construction. Weak readout generalization and imperfect intervention fidelity limit conclusions about the representation independently of the controller. In the secondary analysis, A2 was chosen through outcome-informed pilots, and later design comparisons reused the final math items. Those comparisons are retrospective. The 256-item confirmation tests fixed vectors on fresh project items, not a universal selection rule. Its behavioral audit measures directional alignment and magnitude, not mediation or exact profile matching. Its 95.1\% elicited baseline also limits headroom for improvement. Transfer cells share baselines, and their intervals are exploratory. Equal vector norm does not match achieved behavioral dose. Joint bootstrap and disjoint-baseline sensitivities address uncertainty in average-shift similarity, but duration remains consequential; conditioning on responses that reach a window selects a post-treatment population. Behavioral changes alongside outcomes do not identify mediation.

\paragraph{Dynamic mechanisms.}
The present analyses do not identify the dynamics underlying steering failures. One hypothesis is that some interventions induce repetitive dynamics even when activations remain statistically familiar. A subsequent study could test recurrence, responses to brief perturbations, and recovery of productive output, controlling intervention dose and decoding. Static activation density or projection onto a domain axis would not by itself establish an attractor or its causal role.

\paragraph{Reward learning and evaluation.}
Each historical RL initialization has one training seed. The earlier grader and a temperature-variable defect in the inspected source limit interpretation of the learning process; absent rollout logs and execution hashes prevent reconstructing their exact effects. This material is confined to a secondary appendix analysis. Corrected evaluation establishes outcomes for saved vectors, not a repaired training procedure. The historical forced-answer protocol uses unequal additional computation. The matched-total-budget control is retrospective and uses a fixed answer cue and direct forced-answer parser; it does not equate wall-clock or prefill costs. The clean replication repairs the documented defects and reports three seeds per initialization, but reuses historical evaluation items and tests only natural completion at one budget. Its intervals condition on trained vectors; new-vector annotation, matched-budget elicitation, and comparison with other learned initializations remain untested. Open-task preferences remain judge-dependent and conditional on paired answer availability; neither judge agreement nor zero filling resolves that selection issue.

\section*{Ethical Considerations}
The atlas includes harmful-request and hazardous-knowledge benchmarks. Higher WMDP-Cyber accuracy measures greater hazardous knowledge, not safer behavior, and the present interventions are not validated as safety improvements. We report aggregate patterns without reproducing harmful instructions. Moral and creative quality judgments are limited operationalizations and should not be treated as objective measures of ethical worth or scientific value. Steering methods that alter reasoning or completion may also affect safety in ways not captured by the evaluated objectives. Any release should preserve source-dataset access conditions and licenses and avoid exposing identifying or harmful trace content unnecessarily.
